\begin{solapartat}
\punts{0,50} Si imposem la conservació del nombre de nucleons i de la càrrega elèctrica, tenim:
\[ {}^{210}_{84}Po \rightarrow {}^{4}_{2}He + {}^{206}_{82}Pb \]
\punts{0,25} Per trobar l'activitat després d'una setmana utilitzem l'equació de l'activitat $A(t) = A_0\,e^{-\lambda t}$, on l'activitat inicial: $A_0 = 1,66\cdot 10^{14}\,\frac{Bq}{g}\cdot 5\cdot 10^{-3}\ g = 8,3\cdot 10^{11}\ Bq$,

\punts{0,25} Necessitem el coeficient de desintegració $\lambda$, que obtenim del temps de semidesintegració:\\
$A(t) = \frac{A_0}{2} = A_0\,e^{-\lambda t_{1/2}}$. Per tant: $t_{1/2} = \frac{ln2}{\lambda}$.\\
Directament obtenim: $\lambda = \frac{ln2}{t_{1/2}} = \frac{ln2}{138} = 5,023\cdot 10^{-3}\,dies^{-1}$

\punts{0,25} L'activitat al cap de 7 dies: $A = A_0\,e^{-\lambda t} = 8,3\cdot 10^{11}\,e^{-5,023\cdot 10^{-3}\cdot 7} = 8,01\cdot 10^{11}\,Bq$
\end{solapartat}

\begin{solapartat}
\punts{0,25} El treball efectuat pel camp elèctric és menys l'increment d'energia potencial:
\[ W = -\Delta U = -q\cdot\Delta V \]
\punts{0,25} En portar el primer electró des de l'infinit (distància molt gran), el camp elèctric i potencial elèctric existents són els creats per la partícula alfa. El potencial inicial és nul, $V_i = 0$, i el final serà $V_{f1} = \mathrm{k}\frac{\mathrm{q}_\alpha}{\mathrm{r}}$, on $r$ és la distància final, $r = 0,6\times 10^{-10}$~m. Per tant:
\[ \Delta V_1 = 8,99\cdot 10^{9}\cdot\frac{2\cdot 1,602\cdot 10^{-19}}{0,6\cdot 10^{-10}} = 48,006\ V. \]
I el treball efectuat pel camp elèctric durant el desplaçament del primer electró és:
\[ W_{e1} = -\mathrm{q}_e\cdot\Delta V_1 = 1,602\cdot 10^{-19}\cdot 48,006 = 7,69\cdot 10^{-18}\ J \]
\punts{0,25} En portar el segon electró des de l'infinit, el potencial elèctric existent és el creat per la partícula alfa i l'electró. El potencial inicial és nul i el final serà: $V_{f2} = \mathrm{k}\frac{\mathrm{q}_\alpha}{\mathrm{r}} + \mathrm{k}\frac{\mathrm{q}_e}{2\mathrm{r}}$.

Per tant:
\[ \Delta V_2 = 8,99\cdot 10^{9}\cdot\frac{2\cdot 1,602\cdot 10^{-19}}{0,6\cdot 10^{-10}} - 8,99\cdot 10^{9}\cdot\frac{1,602\cdot 10^{-19}}{2\cdot 0,6\cdot 10^{-10}} = 36,005\ V. \]
I el treball efectuat pel segon electró és:
\[ W_{e2} = -\mathrm{q}_e\cdot\Delta V_2 = 1,602\cdot 10^{-19}\cdot 36,005 = 5,77\cdot 10^{-18}\ J \]

\textit{L'energia potencial final:}

\punts{0,50} L'energia potencial de la configuració final és menys el treball total efectuat pel camp, ja que l'energia potencial inicial era nul·la:\\
$W_{total} = W_{e1} + W_{e2} = 7,69\cdot 10^{-18} + 5,77\cdot 10^{-18} = 1,346\cdot 10^{-17}\ J$\\
$\Delta U = U_f - U_i = -W$. Per tant, l'energia potencial final: és $U_f = -1,346\cdot 10^{-17}\ J$

\textit{O alternativament:}

\punts{0,50} Calculem l'energia potencial electroestàtica del sistema final com a suma dels termes d'energia potencial per parelles:
\[ U_{Sistema} = U_{e-He} + U_{e-He} + U_{e-e} = 2U_{e-He} + U_{e-e} \]
\begin{align*}
U_{Sistema} &= 2\mathrm{k}\frac{\mathrm{q}_\alpha\mathrm{q}_e}{\mathrm{r}} + \mathrm{k}\frac{\mathrm{q}_e\mathrm{q}_e}{2\mathrm{r}} = 2\cdot 8,99\cdot 10^{9}\cdot\frac{2\cdot 1,602\cdot 10^{-19}\cdot(-1,602\cdot 10^{-19})}{0,6\cdot 10^{-10}}\\
&\quad + 8,99\cdot 10^{9}\cdot\frac{(-1,602\cdot 10^{-19})^2}{2\cdot 0,6\cdot 10^{-10}} = -1,346\cdot 10^{-17}\ J
\end{align*}
\end{solapartat}
